\SourceLeaf{150}advantages together. The\,method\,of\,least\ squares is absolutely the most advantageous in every case and for every law of errors. The comparison of the precision of the results with that of the observations, which in my Theoria I based on the law of errors $e^{-hhxx}$, remains generally valid. One also has the advantage that everything is proved and worked out through very clear and simple analytical developments. This hardly seems to be the case with \textsc{Laplace}'s principle and treatment. In particular, the extension of his conclusion from 2 unknown quantities to any number does not yet seem sufficiently evident.

\noindent . . . . . . . . . . .



\Needspace{6\baselineskip}
\begin{center}
\textsc{Gauss} to \textsc{Schumacher}. Göttingen, 2 February 1825.
\end{center}

\noindent . . . . . . . . . . I would in fact like to dispense altogether with the so-called probable errors, since they depend on a hypothesis. One may nevertheless calculate them by multiplying the mean errors by 0,6744897.

\noindent . . . . . . . . . . .



\Needspace{6\baselineskip}
\begin{center}
\textsc{Gauss} to \textsc{Olbers}. Göttingen, 15 March 1827.
\end{center}

\noindent I am much obliged to you, my dearest friend, for kindly sending the three issues of the Philosophical Magazine. I have looked through \textsc{Ivory}'s articles, and I would have to say that I was quite taken aback by them had I not already been prepared by the paper in the Philosophical Transactions (1824), which I mentioned in my last letter. I had long esteemed Mr I.\ as an acute mathematician with great skill in calculation. I still stand by \textls{this} judgment. On the main point of the paper in question, however, I find the logical\space%
